% 第 4 章 与现有工作对比及应用前景
% 本章引用键（实际使用）：tvla, dudect, rtlf, ctgrind, microwalk, kyberslash, ctverif, binsec, clangover, azouaoui2021, jancar2022
% 表号接第 3 章（表 6–11），本章从表 12 起；本章不放图
\setcounter{table}{11}

\section{与现有工作对比及应用前景}

\subsection{泄漏检测与评估方法对比}

软件时间侧信道的评估分为三条路线：统计检测、动态分析与形式化验证。本作品不提出新的检测原语，贡献在于把三条路线组织进同一个以产物为中心、可审计的流程。第 2 章讲这套流程怎么搭，本章讲它与既有路线各自回答什么问题、答案能用到哪一步。

统计检测把程序当黑盒。TVLA\cite{tvla} 比较固定类与随机类的轨迹分布，dudect\cite{dudect} 把同一判据搬到在线计时，用最大 $t$ 与上游阈值给二值结论，RTLF\cite{rtlf} 比较分布与分位，并用 bootstrap 控制第一类错误。三者都只回答“两类样本能不能分开”，不回答差异来自哪条指令；类对构造、样本顺序与运行间漂移都会改变结论。第 3.12 节的 M2 是一个现成例子：RTLF 读成两臂都可区分，dudect 的阳性落在已修臂，预注册的均值检验却两个臂都未拒绝。

动态分析在程序内部找秘密依赖。ctgrind 沿用 Valgrind 的秘密标记，报告秘密进入分支条件或访存下标的时刻\cite{ctgrind}；Microwalk 用执行轨迹与互信息找与秘密关联最强的指令\cite{microwalk}；varlat 检查除法一类可变时延指令的操作数是否含秘密\cite{kyberslash}。这类工具能给出函数级甚至指令级的落点，但只覆盖实际执行过的输入，也不回答“这台机器上能不能测到”。定位阳性与计时阳性是两件事，M2 的三处除法正是两者分离的例子（第 3.4 节、第 3.7 节）。

形式化验证在给定模型与边界内给出另一种结论：证明控制流与访存不依赖秘密。ct-verif 把常数时间归结为产品程序的安全性，在中间表示层完成这类证明\cite{ctverif}；Binsec/Rel 用关系符号执行在二进制层做有界检查，并给出反例路径\cite{binsec}。它们的成本高，通常不建模指令时延、缓存与流水线，所以能回答“是否依赖秘密”，不能回答“本机计时是否可区分”。第 3.4 节 M4 的不安全判定与 ct 臂的有界安全、M5 的 1088/1088 路径完成，都属于有界结论。

表~\ref{tab:compare} 把三条路线与本作品的定位放在同一组维度上。每一行的最后一列写出这个差异导致的实践后果，而不只是一个勾。

\begin{table}[H]
\centering
\caption{泄漏检测与评估方法对比}
\label{tab:compare}
{\footnotesize
\begin{tabularx}{\textwidth}{|>{\raggedright\arraybackslash}p{2.7cm}|>{\raggedright\arraybackslash}p{2.3cm}|>{\raggedright\arraybackslash}p{2.3cm}|>{\raggedright\arraybackslash}p{2.5cm}|>{\raggedright\arraybackslash}p{2.3cm}|>{\raggedright\arraybackslash}X|}
\hline
方法 & 本机可测性 & 函数级定位 & 产物级证明 & 未检出的量化 & 多机制/防护复测 \\
\hline
TVLA\cite{tvla}、dudect\cite{dudect} & 是（黑盒计时） & 否 & 否 & 一般不报，阴性无预算 & 需另行组织；后果：换产物要整体重判 \\
\hline
RTLF\cite{rtlf} & 是（第一类错误有界） & 否 & 否 & 随样本量提高功效，不报 MDE & 需另行组织；后果：分布差异无法归因到指令 \\
\hline
ctgrind\cite{ctgrind}、Microwalk\cite{microwalk} & 否（不回答本机可测性） & 是（分支、访存、轨迹） & 否 & 不适用（只报依赖） & 单一机制类；后果：定位阳性易被读成可利用 \\
\hline
varlat\cite{kyberslash} & 否 & 是（除法等可变时延指令） & 否 & 不适用 & 针对单一指令类；后果：只覆盖已知机制，漏掉新形态 \\
\hline
ct-verif\cite{ctverif}、Binsec/Rel\cite{binsec} & 否 & 是（函数或内核） & 是（给定边界内；ct-verif 在中间表示层） & 不适用（二值结论） & 单一函数或内核；后果：不建模指令时延，答不了本机计时 \\
\hline
本作品（A-16、A-24、A-26、A-27） & 是（预注册 + Holm + MDE） & 是（多工具交叉） & 部分（M4、M5 有界） & 报告 MDE 与样本量 & 5 模块统一规程、防护前后同规程复测；后果：结论绑定产物、可按层追溯 \\
\hline
\end{tabularx}
\par\smallskip
{\footnotesize \textbf{表注：}三条路线回答不同问题：统计检测判“是否可区分”，动态分析判“哪里依赖秘密”，形式化判“给定边界内是否常数时间”。本作品不改变这些判据，只把三类证据登记在同一产物上；“部分”指 M4、M5 的有界验证，不是全实现认证。}
}
\end{table}

五列维度各有用途。本机可测性决定结论能否用于发布判定，函数级定位决定阳性能否落到代码，产物级证明决定结论是否覆盖所有输入，未检出的量化决定阴性有多强，多机制与防护复测决定一次评估能否复用。五列合起来，读者才能判断一条结论该信到哪一步。

本作品把三条路线当作组件来用。同一个二进制既进定位工具，也进计时窗口；定位结果、计时判定与防护复测按证据层分开登记；每个阴性附最小可检测效应，每个未执行单列；再用 90 格排查与 9×5 工具矩阵做交叉（第 3.5 节、第 3.12 节）。这样做的代价是采集与工具运行的投入增加，收益是每个结论都能回答“在哪一层、对哪个产物、在什么预算下成立”。

这种组织的收益在证据归档，不在灵敏度。定位工具还是那些工具，统计检验还是那套检验；改变的是证据的归档方式：同一次运行的对象、产物摘要、编译条件与判定结果绑在一起，其他人也能按同样的顺序复查。

与单机制研究的分界在目标。KyberSlash 以秘密除法与密钥恢复为目标，clangover 追编译器引入的分支，Azouaoui 等系统量化 Kyber 解密与重加密各步骤的泄漏并指向选择密文攻击\cite{kyberslash,clangover,azouaoui2021}。三者都把一条机制追到攻击终点；本作品在“可区分”处停止，不做密钥恢复，把工作重点放在检测的可审计性与多机制覆盖上。与功耗、电磁研究的关系是互补：它们观测物理量，本作品只观测时间，不覆盖功耗、电磁与故障注入。

\textbf{定位阳性不等于本机可利用。}M2 的三处除法在定位层稳定出现，varlat 报 1280 errors / 3 contexts，端到端与除法热点微基准都未检出，最小可检测效应 24.7--67.1 ns 远大于除法的可能效应。定位阳性只说明机制存在；是否可利用要由测量回答。% A-18 % A-03 % A-04

\textbf{锁频的价值是确定性而不是降噪。}不锁频与锁频的 CV 比为 0.964、1.006、1.089，方向不一致；锁频提供的是可核对、可恢复的工作点，使“这次采集在什么频率下完成”可被审计。% A-29

\textbf{同一机制在不同优化级别与架构上形态不同。}M1 在 -Os 下被编成按位测试后条件跳转，-O2/-O3 的命中位置更多（8、15 个）却计时未检出；aarch64 上位域提取与条件选择取代了分支，机制不复现。分支更多不一定更易检出，结论必须绑定产物与环境。% A-26 % A-27

表~\ref{tab:compare} 证明三条路线各有自己的问题域，也证明本作品的证据组织方式能把这些答案放在同一张表上；它不证明本作品在灵敏度、完备性或成本上优于所有对照路线。在单机、单编译器、时间单通道的条件下，本作品的覆盖面小于专用形式化工具，定位精度也不如把单一工具做到极致的场景。它的差异在于每个阴性带预算、每个未执行单列、每个结论绑定产物，评审者可以沿着证据层往回查。

\subsection{应用场景与推广边界}

本作品可以作为密码库发布前检查、构建配置回归、防护方案评估与教学评测的本地工具；它不替代独立审计、形式化认证或真实远程评估。四个场景的共同前提是：使用者手里有源码或二进制，能在一台受控机器上重复运行，并愿意记录编译条件与样本量。对 44 名密码库开发者的调查显示，开发者认为时间侧信道的防护并不难实施，落差在工具采用与工程投入\cite{jancar2022}。

后量子密码库发布前检查由库开发者使用。时机是版本冻结、发布候选构建产出之后；输入是待发布的二进制、编译条件与公开向量；输出是定位层的机制清单和固定机器上的计时判定。定位层用反汇编、Memcheck 与 varlat，成本低，可以在每次候选构建上跑；计时层需要锁频与裸金属，按版本或按周在固定机器上跑。

小团队可以只跑定位层，把计时层留给有裸金属的机器；两层的结果分开登记，不互相代替。得到的是“目标机制是否存在、本机是否可区分”，得不到“全库常数时间认证”。第 3.12 节的 9×5 矩阵说明单工具既可能漏掉机制，也可能给出相反读法，发布判定至少要两类工具同向。

编译器与构建配置变更回归由平台维护者使用。时机是升级编译器、改优化级别或换目标架构之前；clangover 表明这些变更会静默改变产物中的分支形态\cite{clangover}。本作品按“二进制摘要 + 优化级别”组织判定，先比对新旧产物的反汇编与 Memcheck，再在同一锁频工作点复测计时。得到的是新旧产物的差异，得不到“新版本安全”的结论：新产物可能引入新机制，也可能只是把旧机制挪了位置。面向 CI 的回归门禁尚未实现，列为后续工作。

防护方案的工程评估由实现者与评审者使用。时机是候选补丁进入主线之前；输入是同一份源码的防护前后两版、同一组公开向量与同一锁频规程；输出是定位结果、计时判定与成对开销。M1 的掩码选择使解封装 +1.51\%，M2 的除法替代使 $-2.92\%$，前者变慢、后者变快，说明防护代价不能用一个方向概括。得到的是“这个方案在这个产物上的代价与效果”，得不到“所有方案中的最优选择”。% A-25

教学与评测把 M3、M4 的人工校准夹具与 M5 阴性对照当作工具测试集。时机是工具版本冻结或课程实验开始之前；教师或工具开发者用已知机制的夹具检查工具能否识别，用阴性对照检查会不会误报；输入是夹具源码与固定输入类，输出是识别、漏报与误报的清单。得到的是“工具能否识别典型机制、会不会误报”，得不到“工具通过了认证”。M5 上运行过的 6 种工具误报 0，分母是 6 个工具 × 1 个对象，只说明本次没有误报。% A-24

三类任务不适合本作品。跨机器的绝对耗时比较不在范围内，因为时间证据不跨环境合并；功耗、电磁与故障注入未覆盖，因为本作品只观测时间；真实远程链路未覆盖，因为只有同机网络场景。把这些任务交给本作品，会得到看似相关、实际越界的结论。

推广边界有五条，每条都直接约束上面的场景，也对应一项未覆盖的内容。结论绑定在一台 x86\_64 裸金属、clang 18.1.3 与所列优化级别上，不外推到其他 CPU 与编译器。评估在本地白盒下进行，不提供密钥恢复，也不评估功耗、电磁与故障注入。aarch64 只是第二架构，不是第二测量机；网络只是同机模拟，不是真实远程。

真实缺陷采用上游补丁，自主修复只落在 M3、M4 两个人工夹具上；KyberSlash 官方修复提交的残留分支只定位、未修复\footnote{该提交号为 272125f，定位与排查见第 3.5 节。}。题面 12 条按 10 条满足 + 2 条部分登记，缺口与原因见附录 A。

本作品不能替代独立审计、形式化认证或真实远程评估，结果也不构成安全性结论；它能做的是把本地可复现的证据整理成可追溯的记录，让发布者、评审者与教学者知道每一条结论的证据层与适用边界。
